% Deutsche Parallelfassung zu abstract.tex

\cleardoublepage
\chapter*{Zusammenfassung}
\phantomsection
\addcontentsline{toc}{chapter}{Zusammenfassung}

Ein Ingenieur ändert die Krümmung eines Eisenbahn-Alignments. Das
Entwurfssystem zeichnet sofort eine neue Kurve, doch die ingenieurtechnischen
Folgen bleiben nicht innerhalb dieser Zeichnung. Messungen beziehen sich auf
ein physisches Gleis; GIS-Daten ordnen es in eine größere räumliche Umgebung
ein; CAD- und Austauschmodelle beschreiben ausgewählte Konstruktionen; Regeln
und Entscheidungen bestimmen, ob die Änderung zulässig ist. Über diese
Repräsentationen hinweg muss das Alignment identifizierbar bleiben, während
beabsichtigte, realisierte und beobachtete Zustände unterscheidbar bleiben
müssen. Das zentrale Problem ist daher nicht allein der Austausch von
Geometrie. Es besteht darin, Alignment-spezifisches Ingenieurwissen in dem
Moment verfügbar zu machen, in dem eine Änderung interpretiert, bewertet oder
entschieden wird.

Der Ausgangspunkt war konkreter. Der Entwurf von Eisenbahnübergängen muss den
Wechsel zwischen Krümmungszuständen gestalten, während ein Fahrzeug physisch
gezwungen ist, seinem Schienenpaar zu folgen. Krümmung, Überhöhung,
Geschwindigkeit und ihre Änderungsraten wirken durch diese Zwangsführung auf
Fahrzeug, Gleis, Komfort, Sicherheit und Verschleiß. Der ursprüngliche Beitrag
\emph{Berlinish} interpretiert bekannte Übergangsformen durch eine
komponierbare Grammatik aus einlaufender Halbwelle, optionalem Klothoidenkern
und auslaufender Halbwelle. \emph{axtranNew} entstand als vorgesehener
Berechnungskern, der solche Strukturen zusammensetzt, löst, vergleicht und
optimiert. Diese Namen bezeichnen einen ursprünglichen Forschungsbeitrag und
eine Implementierungslinie; sie werden hier nicht als genehmigte Begriffe des
Knowledge Kernel eingeführt.

Diese Thesis entwickelt das \emph{Alignment-Based Information Modelling
(AIM)} als wissensbasierte, repräsentationsbewusste und
implementierungsfähige Antwort. Ihre zentrale These lautet, dass ein Alignment
über seine intrinsische konstruktive Struktur modelliert werden soll,
insbesondere über seine longitudinale Ordnung und seinen Krümmungsverlauf, und
anschließend durch ausdrückliche Beziehungen zwischen Zustand, Realisierung,
Beobachtung und Bewertung zu interpretieren ist. Eine Kurve ist eine
notwendige geometrische Realisierung dieser Struktur, aber nicht das gesamte
Engineering Object.

Die Arbeit leistet eine begriffliche Trennung von Objektidentität,
konstruktiver Definition, metrischer und physischer Realisierung,
Repräsentation, Beobachtung, Wissen und Entscheidung; eine
krümmungsgetriebene Formulierung von Zuständen und Operatoren; eine
Alignment-spezifische ingenieurtechnische Behandlung von Übergängen,
Überhöhung, Nebenbedingungen und Realisierung; eine referenzorientierte
Berechnungspipeline einschließlich hybrider Rekonstruktion und strukturierter
Optimierung; sowie eine Erklärungsfolge, die diese Beiträge vom
Ingenieurproblem bis zur Anwendung miteinander verbindet.

Die weitere ufAIM-Umgebung folgt aus der Notwendigkeit, dieses Übergangswissen
zu bewahren und operativ zu verwenden: Eine Transition-Registry strukturiert
Funktionen und Familien; TransEd macht sie untersuchbar; SPOT trägt Identität
und Zustand konkreter Ingenieurobjekte; alignmentOS macht berechnetes
Alignment-Wissen nutzbar; der Knowledge Kernel stabilisiert Terminologie und
Verantwortung; Research prüft und erweitert Aussagen; und die Thesis erklärt
ihren Zusammenhang. Das resultierende Framework reicht weiter als sein
Ursprung, doch dieser Ursprung begründet seine erste ingenieurwissenschaftliche
Notwendigkeit.

Die mathematischen und numerischen Entwicklungen zeigen, wie Geometrie aus
Krümmung abgeleitet, wie relevante Zustände und Nebenbedingungen bewertet und
wie Alignment-Probleme für die Berechnung formuliert werden können. Sie
implizieren nicht, dass ein Solver Ingenieurwissen oder Entscheidungsautorität
besitzt. AIM bietet vielmehr eine disziplinierte Möglichkeit, Bedeutung bei
wechselnden Repräsentationen zu bewahren und Ingenieurwissen dort anzuwenden,
wo Folgen entstehen.
